\paragraph{Why one word per band, not one per row.} Storing each row in its own 9-bit integer, as
$\mathit{valueToBand}[d-1][b][r]$, would read more simply: a row would be a plain array element, with no shift or mask. It
would not be faster. Taking one row out of the packed word, $(M\gg 9)\,\&\,511$, costs one or two instructions (one on
ARM, which has a bitfield-extract instruction), about the same as loading a separate row from memory, and the shifts in
$P(M)$ are folded into the ORs. Most of the solver's work, though, is on a whole band at once: the band rule's AND with
\textsc{PossibleCells}, the test for an empty mask, the check whether a mask has changed, the stack rule's AND with
$\mathrm{cols}(\mathit{keptColumns}_b)$, the updates when a digit is placed, and the candidate counters of \cref{alg:propagate}, which
handle all 27 cells of a band in each instruction. Each of these is one instruction on a 27-bit word and would be three
on separate rows. Every guess also copies the whole state, and with 16-bit rows \texttt{valueToBand} and \texttt{checked}
would each grow from 108 to 162 bytes. The packed word makes the occasional per-row read slightly less direct so that
every per-band operation is a single instruction.

